% !TEX program = xelatex

\documentclass[svgnames]{amsart}
\usepackage[paperwidth=6in, paperheight=8in, top=20mm, bottom=18mm, left=10mm, right=10mm, head=14pt]{geometry}

\usepackage{amsmath, amsfonts, amssymb, amsthm}
\usepackage{enumitem}
\setlist[enumerate,1]{label=\arabic*., leftmargin=*, itemsep=0.6em}
\setlist[enumerate,2]{label=(\alph*), leftmargin=1.7em, itemsep=0.25em, topsep=0.25em}

\usepackage[T1]{fontenc}
\usepackage[sfdefault, lining, scale=0.9]{FiraSans}
\usepackage{kmath}

\usepackage{xcolor}
\usepackage{scrlayer-scrpage}
\ohead{\color{blue!35!black} \scshape VM}
\cfoot*{\pagemark}

\usepackage{tikz}
\usetikzlibrary{arrows.meta, positioning}

\setlength{\parindent}{0pt}

\newcommand{\comp}[1]{\overline{#1}}
\newcommand{\dist}{\operatorname{d}}
\newcommand{\ecc}{\operatorname{ecc}}
\newcommand{\diam}{\operatorname{diam}}
\newcommand{\rad}{\operatorname{rad}}

\title{SMS 2101: Graph Theory}
\date{}

\begin{document}
\maketitle

\begin{center}
\small
Unless stated otherwise, all graphs are finite, simple, and undirected. \\
For a graph $G$, take $n=|V(G)|$ and $m=|E(G)|$.
\end{center}

\section*{Foundations}

\begin{enumerate}
\item Let $G$ have vertex set $\{1,2,\ldots,7\}$ and edge set
\[
 E(G)=\{12,13,14,23,25,35,36,45,47,56,67\}.
\]
\begin{enumerate}
  \item Find the order, size, degree of each vertex, and degree sequence of $G$.
  \item Let $H$ have vertex set $\{1,2,3,5\}$ and edge set
  $\{12,13,23,25\}$. Is $H$ a subgraph of $G$? Is it induced?
  \item Find the degree sequence and size of $\comp G$ without listing all
  its edges.
\end{enumerate}

\item Decide whether each description is possible. Justify each answer.
\begin{enumerate}
  \item At a meeting of eleven people, each person shakes hands with
  exactly five others.
  \item A graph has degree sequence $(6,5,4,4,3,2,2,2)$.
  \item A graph of order $9$ has one vertex of each degree
  $0,1,\ldots,8$.
  \item There exists a graph of order $n\geq 2$ in which all vertices have different degrees.
\end{enumerate}

\item For the graph $G$ in Problem~1, determine whether each sequence is
a walk, trail, path, closed walk, or cycle, and find its length.
\[
\begin{array}{lll}
W_1=(1,2,3,5,2,1),&
W_2=(7,4,5,3,6,7),&
W_3=(2,1,3,2,5,6,3).
\end{array}
\]

\item Apply the Havel-Hakimi algorithm to each vector below. State
whether it is graphical. If it is, construct a graph having that degree
sequence.
\[
 (5,4,4,3,3,3,2),
 \qquad
 (4,4,3,1,1,1).
\]
Record every reduced vector used by the algorithm.

\item For the graph $G$ in Problem~1:
\begin{enumerate}
  \item find all geodesics from vertex $2$ to vertex $7$;
  \item find the eccentricity of every vertex; and
  \item determine $\rad(G)$, $\diam(G)$, and the centre of $G$.
\end{enumerate}

\item Let
\[
\begin{aligned}
E(G_1)&=\{12,23,34,45,51,13\},\\
E(G_2)&=\{ce,ea,ad,db,bc,ca\}.
\end{aligned}
\]
\begin{enumerate}
  \item Give an isomorphism from $G_1$ to $G_2$.
  \item The graphs $C_6$ and $2K_3$ have the same degree sequence.
  Explain why they are not isomorphic.
\end{enumerate}

\item Let $Q$ have vertex set
\[
 \{a,b,c,x,y,z,w\}
\]
and edge set
\[
 \{ax,ay,by,bz,cx,cz,cw\}.
\]
\begin{enumerate}
  \item Find a bipartition of $Q$.
  \item Is this bipartition unique, apart from interchanging its two parts?
  \item Add the edge $ac$. Exhibit an odd cycle in the resulting graph.
\end{enumerate}

\item For each statement below, decide whether the information guarantees
that $G$ is a tree. Prove your answer or give a counterexample.
\begin{enumerate}
  \item $G$ is connected, has $12$ vertices, and has $11$ edges.
  \item $G$ is acyclic, has $12$ vertices, and has $11$ edges.
  \item $G$ has $12$ vertices and $11$ edges.
\end{enumerate}

\item Let $T$ be the tree with
\[
V(T)=\{a,b,c,d,e,f,g,h,i\},\qquad
E(T)=\{ab,bc,cd,de,cf,cg,gh,hi\}.
\]
Find its pendant vertices, diameter, radius, and centre. Also show the
successive stages of the leaf-deletion procedure used to find the centre.

\end{enumerate}

\section*{Core techniques}

\begin{enumerate}[resume]

\item For each of $K_5$, $K_{2,4}$, and $K_{3,3}$, decide whether it is
Eulerian and whether it is Hamiltonian. Whenever the answer is yes,
exhibit an Eulerian circuit or a Hamiltonian cycle. Whenever it is no,
give a reason.

\item Let $R$ have vertices $v_1,\ldots,v_5$ and edges
\[
e_1=v_1v_2,\ e_2=v_1v_3,\ e_3=v_2v_3,\ e_4=v_2v_4,\
e_5=v_3v_5,\ e_6=v_4v_5.
\]
\begin{enumerate}
  \item Write the adjacency matrix $A(R)$ in the vertex order
  $v_1,\ldots,v_5$.
  \item Write the vertex-edge incidence matrix $B(R)$ using the edge
  order $e_1,\ldots,e_6$.
  \item Verify directly that $BB^{T}=A+D$, where $D$ is the diagonal
  matrix of vertex degrees.
\end{enumerate}

\item Use powers of the adjacency matrix from Problem~11 to find:
\begin{enumerate}
  \item the number of walks of length $2$ from $v_1$ to $v_5$;
  \item the number of closed walks of length $3$ starting at $v_1$; and
  \item the total number of triangles in $R$.
\end{enumerate}
Explain how each answer is read from the appropriate matrix entry or
trace.

\item Apply Dijkstra's algorithm to the directed weighted graph below, starting at
$s$. At every iteration, record the permanent vertex, all tentative
distances, and the predecessor of each unsettled vertex. Give the
shortest-path tree and the shortest path from $s$ to $t$.

\begin{center}
\begin{tikzpicture}[
  scale=0.82,
  every node/.style={transform shape},
  vertex/.style={circle, draw=blue!35!black, fill=blue!5,
    minimum size=6mm, inner sep=0pt, font=\bfseries},
  arc/.style={-{Stealth[length=2mm]}, draw=blue!35!black},
  edge label/.style={fill=white, inner sep=1pt, font=\small}
]
\node[vertex] (s) at (0,1) {$s$};
\node[vertex] (a) at (2.4,2.3) {$a$};
\node[vertex] (b) at (2.4,-0.3) {$b$};
\node[vertex] (c) at (5.2,2.3) {$c$};
\node[vertex] (d) at (5.2,-0.3) {$d$};
\node[vertex] (t) at (7.6,1) {$t$};

\draw[arc] (s)--node[edge label, above] {7} (a);
\draw[arc] (s)--node[edge label, below] {2} (b);
\draw[arc, bend left=15] (a) to node[edge label, right] {6} (b);
\draw[arc, bend left=15] (b) to node[edge label, left] {2} (a);
\draw[arc] (a)--node[edge label, above] {4} (c);
\draw[arc, bend left=16] (b) to node[edge label, below left, pos=0.42] {9} (c);
\draw[arc, bend left=16] (c) to node[edge label, above right, pos=0.58] {4} (b);
\draw[arc] (a)--node[edge label, right, pos=0.42] {6} (d);
\draw[arc, bend left=15] (b) to node[edge label, below] {12} (d);
\draw[arc, bend left=15] (d) to node[edge label, above] {5} (b);
\draw[arc] (c)--node[edge label, right] {1} (d);
\draw[arc, bend left=15] (c) to node[edge label, above] {10} (t);
\draw[arc, bend left=15] (t) to node[edge label, below] {6} (c);
\draw[arc] (d)--node[edge label, below] {3} (t);
\end{tikzpicture}
\end{center}

How do the shortest path and distance from $s$ to $t$ change if the
weight of the arc $c\to t$ is changed from $10$ to $2$?

\item Let $G$ be a graph of order $n$.
\begin{enumerate}
  \item Express $\deg_{\comp G}(v)$ in terms of $\deg_G(v)$.
  \item Prove that $A(\comp G)=J-I-A(G)$, where $J$ is the all-ones matrix of order $n$ and $I$ is the identity matrix of order $n$.
  \item If $G$ and $\comp G$ are both $k$-regular, determine $n$ in terms of $k$, and show that $k$ must be even.
\end{enumerate}

\end{enumerate}

\section*{Proof and synthesis}

\begin{enumerate}[resume]

\item Prove the Handshaking Lemma. Deduce that:
\begin{enumerate}
  \item every graph has an even number of vertices of odd degree; and
  \item a $k$-regular graph of odd regularity has even order.
\end{enumerate}

\item Prove that every graph on at least two vertices has two vertices
of the same degree.

\item Prove that every graph on six vertices contains a triangle either
in the graph or in its complement. Show that five vertices do not suffice
by considering $C_5$. Thus establish that $R(3,3)=6$.

\item Prove that if $G$ is disconnected, then $\comp G$ is connected.
In fact, show that any two vertices of $\comp G$ are at distance at most
$2$.

\item Let $G$ be connected. Prove that if $\diam(G)\geq 3$, then
$\diam(\comp G)\leq 3$.

\item Suppose $G$ has order $n\geq2$ and
\[
\delta(G)\geq \frac{n-1}{2}.
\]
Prove that $G$ is connected. Then prove the stronger conclusion
$\diam(G)\leq2$ (and hence, in particular, $\diam(G)\leq3$).

\item Prove that a graph is bipartite if and only if it contains no odd
cycle. For the converse, fix a vertex in each component and partition
the component according to whether distance from that vertex is even or
odd.

\item Prove that a graph $T$ is a tree if and only if there is a unique
path between every pair of its vertices. Use this to prove that every
tree of order $n$ has exactly $n-1$ edges.

\item Prove that every non-trivial tree has at least two pendant vertices.
Then justify the leaf-deletion method: deleting all pendant vertices from
a tree with at least three vertices produces a smaller tree with the same
centre. Deduce that the centre of a tree consists of either one vertex or
two adjacent vertices.

\item
\begin{enumerate}
  \item Prove that a self-complementary graph has order $4r$ or $4r+1$
  for some non-negative integer $r$.
  \item Verify explicitly that $P_4$ is self-complementary.
  \item Let $G$ be self-complementary. Form $H$ from the disjoint union
  of $G$ and a path $a-b-c-d$ by joining every vertex of $G$ to $b$ and
  $c$. Prove that $H$ is self-complementary. Use this construction to
  give a self-complementary graph of order $8$.
\end{enumerate}

\item Prove that every vertex of an Eulerian graph has even degree. Give
an example of a graph in which every vertex has even degree but which is
not Eulerian, and explain which part of the definition fails.
\end{enumerate}

\end{document}
